State Bayes’ rule and the law of total probability.

Law of Total Probability:
The probability in the denominator can be calculated as
\[
\mathbb{P}[B] = \mathbb{P}[B|A]\mathbb{P}[A] + \mathbb{P}[B|A^c]\mathbb{P}[A^c].
\]

Bayes' Rule:
The conditional probability \(\mathbb{P}[A|B]\) is given by
\[
  \mathbb{P}[A|B] = \frac{\mathbb{P}[B|A]\mathbb{P}[A]}{\mathbb{P}[B]} = \frac{\mathbb{P}[A \cap B]}{\mathbb{P}[B]}
\]

Interpretation:
If we know that the event \(B\) occurs, it changes our belief on the probability of \(A\).
Put into classical terms, Posterior = Likelihood × Prior ÷ Evidence.


Describe the approach of Bayesian statistics. Explain the notion of prior, posterior
and conditional distribution. State the updating formula for pmf and pdf. Discuss
the proportionality trick.

Prior Distribution: 
This represents the initial beliefs about a parameter \(\vartheta\) before observing any new data. 
It is denoted by \(f_\vartheta(\vartheta)\) and expresses the uncertainty about the parameter based on prior knowledge.

Conditional Distribution: Given data \(\mathbf{x}\), the likelihood or conditional distribution \(f_{X|\vartheta}(\mathbf{x}|\vartheta)\) quantifies the probability of observing the data under various parameter values.

 Posterior Distribution: After observing data, the prior is updated to a posterior distribution \(f_{\vartheta|X}(\vartheta|\mathbf{x})\), incorporating both the prior beliefs and the likelihood from the observed data.

Updating Formula
For a probability mass function (discrete case):
\[
f_{\vartheta|X}(\vartheta|\mathbf{x}) = \frac{f_{X|\vartheta}(\mathbf{x}|\vartheta) \cdot f_\vartheta(\vartheta)}{f_X(\mathbf{x})}
\]
For a probability density function (continuous case):
\[
f_{\vartheta|X}(\vartheta|\mathbf{x}) = \frac{f_{X|\vartheta}(\mathbf{x}|\vartheta) \cdot f_\vartheta(\vartheta)}{\int f_{X|\vartheta}(\mathbf{x}|\vartheta) \cdot f_\vartheta(\vartheta) \, d\vartheta}
\]

Proportionality Trick
When computing the posterior distribution, the denominator \(f_X(\mathbf{x})\) or \(\int f_{X|\vartheta}(\mathbf{x}|\vartheta) \cdot f_\vartheta(\vartheta) \, d\vartheta\) can be complex. However, since it is a normalizing constant, the posterior is often expressed up to a constant of proportionality:
\[
f_{\vartheta|X}(\vartheta|\mathbf{x}) \propto f_{X|\vartheta}(\mathbf{x}|\vartheta) \cdot f_\vartheta(\vartheta)
\].


Define Bayes risk and Bayesian point estimators. Discuss examples.

Definition:
We consider some probability space \((\Omega, \mathcal{F}, \mathbb{P})\) on which all random variables are defined:
\(\theta \in \mathbb{R}\) has the known prior distribution \(f_\theta\).
The \(\mathcal{X}\)-valued random data \(X\) with \(\theta\)-conditional distribution \(f_{X|\theta}\) are distributed according to
\[
f_X(x) = \int_{\vartheta \in \Theta} f_{X|\vartheta}(x|\vartheta)f_\theta(\vartheta) \, d\vartheta
\]
where \(d\vartheta\) is the Lebesgue measure in the continuous case and the counting measure in the discrete case.
  
A mapping \(T : \mathcal{X} \to \mathbb{R}\) is called (point) estimator for \(\theta\).
For an admissible estimator (measurability, integrability), we assume the existence of a loss function \(L(\theta, T)\).

The function
\[
\rho^L(T) := \mathbb{E}[L(\theta, T(X))]
\]
is called the Bayes risk of \(T\) (under \(L\)).

An admissible estimator \(T^*\) is called Bayes estimator for \(\theta\) (under \(L\)), 
if it minimizes the Bayes risk \(\rho^L(T)\) over all admissible estimators \(T\).

The most common choice of a loss function to determine Bayesian point estimators is the squared error (SE):

\[
L^{\text{SE}}(\theta, T) = (\theta - T(X))^2.
\]

Form of the Bayes estimator under SE:
The Bayes estimator under SE \(T^{*,\text{SE}} : \mathcal{X} \to \mathbb{R}\) is given by
\[
T^{*,\text{SE}}(x) = \mathbb{E}[\theta|X = x],
\]
i.e., by the conditional expectation of \(\theta\) given \(X = x\).
This formula immediately results from the properties and calculation rules of the conditional expectation.
The Bayes estimator can be calculated using the posterior distribution \(f_{\vartheta|X}(\vartheta|x)\):
\[
T^{*,\text{SE}}(x) = \mathbb{E}[\theta|X = x] = \int_\Theta \vartheta \, f_{\vartheta|X}(\vartheta|x) \, d\vartheta.
\]

The Bayes estimator under SE corresponds to the mean of the posterior distribution.
Example
\(\theta|X = x\) possesses a beta distribution with parameters \(x + 1\) and \(n - x + 1\). 
Hence, the Bayes estimator under SE for the one-year claim occurrence probability is given by
\[
T^{*,\text{SE}}(x) = \mathbb{E}[\theta|X = x] = \frac{x+1}{n+2}
\]
if claims have occurred in \(x\) out of the \(n\) observed years.

Introduce the Binomial sampling model with potential applications.

The one-parameter model described by the conditional probability mass function
\[
f_{X|\theta}(x|\vartheta) = \binom{n}{x} \vartheta^x (1 - \vartheta)^{n-x},
\]

i.e., \(X|\theta = \vartheta \sim \text{Bin}(n, \vartheta)\) is called the binomial model.

The binomial model appears in various applications:
\(X\): number of claims caused by a driver with unknown claim occurrence probability per year \(\vartheta\) in \(n\) years.
Happiness data 

The data in the sketched application examples possesses the following unifying structure:
The sample size is fixed and equal to \(n\).
For each outcome in the sample \(i \in \{1, \ldots, n\}\), a certain event of interest may either occur, \(X_i = 1\), or not occur, \(X_i = 0\), each with the same unknown occurrence probability \(\vartheta\).
Outcomes in the sample are conditionally independent given \(\theta = \vartheta\).
Interest is in the number of occurring events in the sample \(X = \sum_{i=1}^{n} X_i\).

Mathematically, this leads to the assumptions
\(X_1, \ldots, X_n | \theta = \vartheta\) are i.i.d. \(\text{Ber}(\vartheta)\) distributed random variables.
\(X = \sum_{i=1}^{n} X_i | \theta = \vartheta \sim \text{Bin}(n, \vartheta)\).


Define the notion of priors conjugate for a sampling model. Discuss this for the
Binomial sampling model. Derive the posterior for a Beta prior.

A class \( \mathcal{P} \) of prior distributions for \( \theta \) is called conjugate for a sampling model \( f_{X|\theta} \) if
\( f_{\theta} \in \mathcal{P} \implies f_{\theta|X} \in \mathcal{P} \)

Binomial Model: Posterior under Beta Prior

\[
\begin{array}{ll}
    \bullet & \theta \sim \text{Beta}(a, b), \; a, b > 0 \\
    \bullet & X|\theta = \vartheta \sim \text{Bin}(n, \vartheta), \\
\end{array}
\]
then \(\theta | X = x \sim \text{Beta}(x + a, n - x + b)\).
Proof:
\[
f_{\theta|x}(\vartheta|x) = \frac{f_{X|\theta}(x|\vartheta) f_{\theta}(\vartheta)}{f_{X}(x)} 
= \frac{1}{f_{X}(x)} \cdot \frac{1}{\beta(a, b)} \vartheta^{a-1} (1 - \vartheta)^{b-1} \cdot \mathbb{1}_{[0,1]}(\vartheta) 
\left( \binom{n}{x} \vartheta^x (1 - \vartheta)^{n-x} \right)
\]
\[
= c(n, x, a, b) \cdot \vartheta^{x+a-1} (1 - \vartheta)^{b+n-x-1} \cdot \mathbb{1}_{[0,1]}(\vartheta) 
\stackrel{(*)}{=} f_{\text{Beta}(a+x,b+n-x)}(\vartheta)
\]

Proportionality trick for step \((*)\):
Our posterior density and the beta density are proportional (only difference is the constant \(c(n, x, a, b)\)), i.e., they share the same shape.
But since both the beta density and our posterior density integrate to 1, they also share the same scale.


Elaborate on sufficient statistics for estimating parameters in the Bayesian setting.

Observing the realized data \(x_1, \ldots, x_n\), Bayes' rule yields:
\[
f_{\vartheta | X_1, \ldots, X_n}(\vartheta | x_1, \ldots, x_n) = \frac{f_{X_1, \ldots, X_n | \vartheta}(x_1, \ldots, x_n | \vartheta) \cdot f_{\vartheta}(\vartheta)}{f_{X_1, \ldots, X_n}(x_1, \ldots, x_n)}
\]
\[
= \vartheta^{\sum_{i=1}^{n} x_i} \cdot (1 - \vartheta)^{n-\sum_{i=1}^{n} x_i} \cdot f_{\vartheta}(\vartheta) \cdot \frac{1}{f_{X_1, \ldots, X_n}(x_1, \ldots, x_n)}
\]

If we now compare the relative conditional probabilities evaluated at \(\vartheta_a\) and \(\vartheta_b\), we see that
\[
\frac{f_{\vartheta|X_1, \ldots, X_n}(\vartheta_a | x_1, \ldots, x_n)}{f_{\vartheta|X_1, \ldots, X_n}(\vartheta_b | x_1, \ldots, x_n)}
= \left(\frac{\vartheta_a}{\vartheta_b}\right)^{\sum_{i=1}^{n} x_i} \cdot \left(\frac{1 - \vartheta_a}{1 - \vartheta_b}\right)^{n-\sum_{i=1}^{n} x_i} \cdot \frac{f_{\vartheta}(\vartheta_a)}{f_{\vartheta}(\vartheta_b)}
\]

i.e., the ratio of the conditional densities/probabilities at the two arguments \(\vartheta_a\) and \(\vartheta_b\) depends on the data 
\(x_1, \ldots, x_n\) only through their sum \(\sum_{i=1}^{n} x_i\).

One can show that
\[
\mathbb{P}(\theta \in A \mid X_1 = x_1, \ldots, X_n = x_n) = \mathbb{P}\left( \theta \in A \mid \sum_{i=1}^{n} X_i = \sum_{i=1}^{n} x_i \right)
\]
for any event \( A \in \mathcal{B}([0, 1]) \).

This means that \( \sum_{i=1}^{n} X_i \) contains all information about \( \theta \) available in the data.
We say that \( \sum_{i=1}^{n} X_i \) is a sufficient statistic for \( \theta \) and \( f_{\theta \mid x_1, \ldots, x_n} \).


Define the notation of predictive distribution. Compute it for the Binomial sampling model with beta prior.

An important feature of Bayesian inference is the existence of a predictive distribution for new observations.
Let \(x_1, \ldots, x_n\) be the outcomes of the trials underlying the chosen sampling model and let \(X_{n+1}\) be 
the random outcome of an additional trial that has not yet been carried out.

Definition 6.2.3 (Predictive Distribution) 

The predictive distribution of \(X_{n+1}\) is the conditional distribution of \(X_{n+1}\) given \(\{X_1 = x_1, \ldots, X_n = x_n\}\).

Application to the Binomial Model:
For conditionally i.i.d. Bernoulli random variables, it holds that

\[
\mathbb{P}(X_{n+1} = 1 | X_1 = x_1, \ldots, X_n = x_n) = f_{X_{n+1} | X_1, \ldots, X_n}(1 | x_1, \ldots, x_n)
\]
\[
= \int f_{X_{n+1}| \theta, X_1, \ldots, X_n}(1 | \vartheta, x_1, \ldots, x_n) f_{\theta | X_1, \ldots, X_n}(\vartheta | x_1, \ldots, x_n) d\vartheta
\]
\[
= \int \vartheta \cdot f_{\theta | X_1, \ldots, X_n}(\vartheta | x_1, \ldots, x_n) d\vartheta
\]
\[
= \mathbb{E}[\theta | X_1 = x_1, \ldots, X_n = x_n] = \frac{a + \sum_{i=1}^{n} x_i}{a + b + n}
\]
\[
\mathbb{P}(X_{n+1} = 0 | X_1 = x_1, \ldots, X_n = x_n) = 1 - \mathbb{E}[\theta | X_1 = x_1, \ldots, X_n = x_n] = \frac{b + \sum_{i=1}^{n} (1 - x_i)}{a + b + n}
\]

What is p-Bayesian coverage? What is frequentist coverage? How can confidence
intervals be computed on the basis of quantiles? What are HPD regions, and what
is their advantage?

A random interval \( \left[ l(\mathbb{X}), u(\mathbb{X}) \right] \)  is said to possess a \( p \)-frequentist coverage for \( \theta \) if
\[
\forall \vartheta : \quad \mathbb{P}(l(\mathbb{X}) < \vartheta < u(\mathbb{X}) | \theta = \vartheta) \geq p.
\]

Interpretation: The standard frequentist confidence interval describes the \(\theta\)-conditional probability that the interval 
will cover the true value \( \vartheta \) before the data is observed.
Frequentist coverage Pre-experimental coverage and prior distribution is irrelevant.
Bayesian coverage Post-experimental coverage, but result depends on choice of prior distribution.

Goal: We will construct confidence intervals with \( p \)-Bayesian coverage that asymptotically also possess p-frequentist coverage.

Idea: Use posterior quantiles to determine a confidence interval
Steps: In order to create a \(1 - \alpha\) quantile-based confidence interval:
Find numbers \(\theta_{\alpha/2} < \theta_{1-\alpha/2}\) such that
\(\mathbb{P}(\theta < \theta_{\alpha/2} \mid X = x) = \frac{\alpha}{2}\),
\(\mathbb{P}(\theta > \theta_{1-\alpha/2} \mid X = x) = \frac{\alpha}{2}\).

The numbers \(\theta_{\alpha/2}, \theta_{1-\alpha/2}\) are the \(\alpha/2\) and \(1-\alpha/2\) posterior quantiles of \(\theta\) given \(X = x\).

If there are \( \theta \)-values outside of the quantile-based interval that possess a higher posterior probability
(density) than some points inside the interval, another approach is required.

A \( (1 - \alpha) \)-HPD region is defined as a subset of the parameter space \( s(x) \subseteq \Theta \) such that

\( \mathbb{P}(\theta \in s(x) \mid X = x) = 1 - \alpha \),
If \( \vartheta_a \in s(x) \) and \( \vartheta_b \notin s(x) \), then 
\( f_{\theta \mid X}(\vartheta_a \mid x) > f_{\theta \mid X}(\vartheta_b \mid x) \).

All points in an HPD region have a higher posterior density than points outside the region.
Note, however, that if the posterior density is multimodal, an HPD region might not be an interval any more!


Introduce the Poisson sampling model with potential applications. Show that
Gamma priors are conjugate for this sampling model.

The one-parameter model described by the conditional probability mass function for the conditionally independent samples
\[
f_{X_i|\theta}(x_i|\vartheta) = \frac{\vartheta^{x_i}}{x_i!} e^{-\vartheta}, \quad \text{for } x_i \in \{0, 1, 2, \ldots\}
\]

i.e., \(X_i|\theta = \vartheta \sim \text{Poi}(\vartheta)\), \(i = 1, \ldots, n\), (i.e., \(X_i|\theta = \vartheta\), 
i.i.d. Poisson-distributed) is called the Poisson model.

If \(\theta \sim \Gamma(a, b), \, a, b > 0\)
\(X_1, \ldots, X_n | \theta = \vartheta \sim \text{Poi}(\vartheta)\)
then \(\left( \theta \big| \{X_1 = x_1, \ldots, X_n = x_n\} \right) \sim \Gamma\left(a + \sum_{i=1}^n x_i, b + n\right)\).

Proof: In the considered situation, it holds that

\begin{align*}
f_{\theta|X_1,\ldots,X_n}(\vartheta|x_1,\ldots,x_n) &= c(x_1,\ldots,x_n) \cdot e^{-n\vartheta} \vartheta^{\sum_{i=1}^n x_i} \cdot f_{\theta}(\vartheta) \\
&= c(x_1,\ldots,x_n) \cdot e^{-n\vartheta} \vartheta^{\sum_{i=1}^n x_i} \cdot \frac{b^a}{\Gamma(a)} \vartheta^{a-1} e^{-b\vartheta} \\
&= \bar{c}(x_1,\ldots,x_n,a,b) \cdot \vartheta^{\sum_{i=1}^n x_i} e^{-n\vartheta} \cdot \vartheta^{a-1} e^{-b\vartheta} \\
&= \bar{c}(x_1,\ldots,x_n,a,b) \cdot \vartheta^{a + \sum_{i=1}^n x_i-1} e^{-(b+n)\vartheta},
\end{align*}
which proves the statement.

Using the properties of the gamma distribution, we find that the posterior expectation of \( \theta \) can,
analogous to the binomial model, be written as a convex combination of prior expectation and sample average:

\[
\mathbb{E}[\theta | X_1 = x_1, \ldots, X_n = x_n] = \frac{a + \sum_{i=1}^n x_i}{b + n}
\]
\[
= \frac{b}{b+n} \cdot \underbrace{\frac{a}{b}}_{\text{prior expectation}} + \frac{n}{b+n} \cdot \underbrace{\frac{\sum_{i=1}^n x_i}{n}}_{\text{sample average}}
\]

Interpretation of the parameters in terms of prior information:
\( b \) is interpreted as the number of prior observations
\( a \) is interpreted as the sum of counts from the \( b \)

Implication for large sample size: Analogous to the binomial model, for large sample size \( n \) , the
information from the data dominates the prior information.

What are exponential families? If they are sampling models, what are convenient
choices of conjugate priors? Explain that the Bernoulli sampling model and the
Poisson model are special cases.

A one-parameter exponential family distribution with parameter \( \phi \) is characterized through a density/probability mass function of the form
\[
f_{X|\phi}(x|\phi) = h(x) \cdot c(\phi) \cdot e^{\phi t(x)}
\]
where \( t(x) \) is the sufficient statistic of the distribution, \( h \) is an arbitrary measurable function, 
and \( c \) is the normalizing constant.


One can show that for a particular prior distribution of the form
\[
f_{\phi}(\varphi) = \kappa(n_0, t_0) \cdot c(\varphi)^{n_0} \cdot e^{n_0 t_0 \varphi}
\]
a one-parameter exponential family model will yield the posterior distribution
\[
f_{\phi|X_1,\ldots,X_n}(\varphi|X_1,\ldots,X_n) \propto f_{\phi}(\varphi) \cdot f_{X_1,\ldots,X_n|\phi}(X_1,\ldots,X_n|\varphi)
\]
\[
\propto c(\varphi)^{n_0+n} \exp \left( \varphi \cdot \left[ n_0 t_0 + \sum_{i=1}^{n} t(x_i) \right] \right)
\]

which is of the same form as the prior distribution with updated parameters \( n_0 + n \) instead of \( n_0 \) and 
\( n_0 \cdot t_0 + n \cdot \bar{t}(x_1,\ldots,x_n) \) instead of \( n_0 \cdot t_0 \) with 
\( \bar{t}(x_1,\ldots,x_n) = \frac{1}{n} \sum_{i=1}^{n} t(x_i) \).
This is the conjugate prior for one-parameter exponential family models
Interpretation
\( n_0 \) can be interpreted as the 'prior sample size',
\( t_0 \) is interpreted as 'prior guess' of \( t(X_i) \) (to be precise, one can show that: \( t_0 = \mathbb{E}[t(X_i)] \)).

E.g. Binomial distribution
For a single binary variable, \( X_1 \), we have:
\(
f_{X_1|\theta}(x_1|\vartheta) = \vartheta^{x_1} \cdot (1 - \vartheta)^{1-x_1} = \left(\frac{\vartheta}{1-\vartheta}\right)^{x_1} \cdot (1 - \vartheta)
\)
\(
= (1 + e^{\varphi})^{-1} \cdot e^{\varphi x_1},
\)
where \( \varphi = \log \left( \frac{\vartheta}{1-\vartheta} \right) \)
We obtain an exponential family model with \( h(x_1) = 1 \), \( c(\varphi) = (1 + e^{\varphi})^{-1} \) and sufficient statistic \( t(x_1) = x_1 \).
The conjugate prior in terms of \( \phi \) is thus given by
\(
f_{\phi}(\varphi) \propto (1 + e^{\varphi})^{-n_0} \cdot e^{n_0 t_0 \varphi}.
\)
Changing variables, this translates into a prior distribution for \( \theta \):
\(
f_{\theta}(\vartheta) \propto \vartheta^{n_0 t_0 -1} (1 - \vartheta)^{n_0(1-t_0)-1},
\)

The Poisson model with parameter \( \vartheta \) can be shown to belong to the class of one-parameter exponential family models with
\( t(x) = x \)
\( \varphi = \log(\vartheta) \)
\( c(\varphi) = \exp(e^{-\varphi}) \).
The conjugate prior in terms of \( \phi \) is thus given by
\(
f_{\phi}(\varphi) \propto \exp(n_0 e^{-\varphi}) e^{n_0 t_0 \varphi}.
\)
This translates into the conjugate prior in terms of \( \theta \) as
\(
f_{\theta}(\vartheta) \propto \vartheta^{n_0 t_0 - 1} e^{-n_0 \vartheta},
\)
the density of a gamma distribution with parameters \( n_0 t_0, n_0 \).
